\subsubsection{Rank Deficiency, Moment Visibility, and Hankel Theory}

We now turn to the structural constraints imposed by the fiber multiplicity distribution on the associated moment matrices and Hankel operators.

\begin{definition}[Moment sequence and truncated Hankel matrix]
\label{def:pom-moment-hankel}
Given the fiber multiplicity function $\nu_m:k\mapsto|\{x\in X_m:d_m(x)=k\}|$, define the \emph{normalized moment sequence}:
\[
\mu_r:=\frac{1}{N_m}\sum_{k}\nu_m(k)\cdot k^r,
\qquad
r=0,1,2,\ldots,
\]
where $N_m=|X_m|=F_{m+2}$. The \emph{truncated Hankel matrix} of order $L$ is:
\[
H_L:=
\begin{pmatrix}
\mu_0 & \mu_1 & \mu_2 & \cdots & \mu_L \\
\mu_1 & \mu_2 & \mu_3 & \cdots & \mu_{L+1} \\
\mu_2 & \mu_3 & \mu_4 & \cdots & \mu_{L+2} \\
\vdots & \vdots & \vdots & \ddots & \vdots \\
\mu_L & \mu_{L+1} & \mu_{L+2} & \cdots & \mu_{2L}
\end{pmatrix}.
\]
\end{definition}

\begin{proposition}[Hankel rank bound from value set complexity]
\label{prop:pom-hankel-rank-bound}
Under the notation of Definition~\ref{def:pom-moment-hankel}, let $\mathcal{V}_m:=\{d_m(x):x\in X_m\}$ denote the value set of the fiber index function. Then, for any $L\ge0$, the Hankel matrix $H_L$ satisfies:
\[
\mathrm{rank}(H_L)
\le
|\mathcal{V}_m|
=
\exp(O(\sqrt m)).
\]
In particular, $H_L$ is \emph{severely rank-deficient} when $L\gg\sqrt m$, with the rank-to-size ratio decaying as:
\[
\frac{\mathrm{rank}(H_L)}{L+1}
=
O\!\left(
\frac{\exp(\sqrt m)}{L}
\right).
\]
\end{proposition}

\begin{proof}
The moment sequence $\{\mu_r\}$ is supported on the finite set $\mathcal{V}_m$, i.e., it can be written as:
\[
\mu_r
=
\frac{1}{N_m}\sum_{k\in\mathcal{V}_m}\nu_m(k)\cdot k^r
=
\sum_{k\in\mathcal{V}_m}w_k\cdot k^r,
\]
where $w_k:=\nu_m(k)/N_m\ge0$ and $\sum_k w_k=1$.

This is precisely the $r$-th moment of a discrete probability measure supported on $\mathcal{V}_m$. By the classical theory of moment problems (see~\cite{Akhiezer1965ClassicalMomentProblem}), the Hankel matrix $H_L$ associated with such a measure has rank at most $|\mathcal{V}_m|$, regardless of the order $L$.

From Remark~\ref{rem:pom-fiber-value-set-complexity}, we have $|\mathcal{V}_m|=\exp(\Theta(\sqrt m))$. Therefore:
\[
\mathrm{rank}(H_L)
\le
|\mathcal{V}_m|
=
\exp(O(\sqrt m)).
\]

When $L$ is chosen to be polynomial in $m$ (e.g., $L=\Theta(m)$), the Hankel matrix $H_L$ is a square matrix of size $(L+1)\times(L+1)=\Theta(m^2)$, while its rank is only $\exp(O(\sqrt m))$. The rank-to-size ratio is:
\[
\frac{\mathrm{rank}(H_L)}{L+1}
\sim
\frac{\exp(c\sqrt m)}{m}
\to0
\quad\text{as }m\to\infty,
\]
demonstrating severe rank deficiency.
\end{proof}

\begin{remark}[Geometric interpretation: moment variety collapse]
\label{rem:pom-hankel-collapse}
The rank deficiency of $H_L$ can be interpreted geometrically as a \emph{moment variety collapse}. In the language of algebraic statistics, the moment sequence $\{\mu_r\}$ defines a point in the infinite-dimensional space of all moment sequences. The Hankel rank bound implies that this point lies on a variety of dimension at most $|\mathcal{V}_m|=\exp(O(\sqrt m))$, which is exponentially smaller than the ambient dimension.

This collapse is a direct consequence of the Fibonacci factorization structure: the fiber index $d_m(x)$ can only take values in a restricted subset of $\mathbb{N}$ (namely, products of Fibonacci numbers), and this restriction propagates through to the moment sequence, forcing it to lie on a low-rank subvariety.

From the perspective of projection ontology, this phenomenon represents a \emph{structural bottleneck}: the type space $X_m$ has exponential size $F_{m+2}$, but the moment structure is constrained by the subexponential value set $\mathcal{V}_m$. This bottleneck is a manifestation of the \emph{projection principle}: high-dimensional combinatorial structures (Zeckendorf types) collapse onto low-dimensional algebraic objects (moment varieties) when projected through the fiber index map $x\mapsto d_m(x)$.
\end{remark}

\begin{proposition}[Moment visibility threshold]
\label{prop:pom-moment-visibility}
Define the \emph{normalized $r$-th central moment}:
\[
\sigma_r:=\frac{1}{N_m}\sum_{k\in\mathcal{V}_m}\nu_m(k)\cdot(k-\mu_1)^r,
\]
where $\mu_1$ is the first moment (mean fiber index). Then, for $r\ge2$, we have:
\[
\sigma_r
=
\Theta\!\left(
\phi^{r\cdot m}
\right),
\]
i.e., the central moments grow exponentially with the moment order and the type space size.

Moreover, there exists a \emph{visibility threshold} $r_*(m)=\Theta(\sqrt m/\log m)$ such that:
\begin{itemize}
\item For $r\le r_*(m)$: the moment $\sigma_r$ is \emph{statistically visible}, i.e., it can be reliably estimated from a polynomial number of samples.
\item For $r>r_*(m)$: the moment $\sigma_r$ is \emph{statistically invisible}, i.e., it requires exponentially many samples to estimate accurately.
\end{itemize}
\end{proposition}

\begin{proof}
From Theorem~\ref{thm:pom-fiber-indset-factorization}, each fiber index value $k\in\mathcal{V}_m$ is a product of Fibonacci numbers:
\[
k=\prod_{i=1}^rF_{\ell_i+2},
\qquad
\sum_{i=1}^r\ell_i\le\left\lceil\frac{m}{2}\right\rceil.
\]
The asymptotic formula $F_n\sim\phi^n/\sqrt5$ implies:
\[
k
\sim
\phi^{\sum_i(\ell_i+2)}/5^{r/2}
\le
\phi^{m/2+2r}.
\]
The maximum fiber index $k_{\max}$ is achieved when the partition $\{\ell_i\}$ is maximally concentrated, yielding:
\[
k_{\max}
\sim
\phi^{m/2+O(1)}.
\]

The first moment $\mu_1$ is the average fiber index:
\[
\mu_1
=
\frac{1}{N_m}\sum_{x\in X_m}d_m(x)
=
\Theta(\phi^{m/2}),
\]
as verified numerically in Figure~\ref{fig:pom-fiber-mean-asymptotic}.

For the $r$-th central moment, the dominant contribution comes from values $k$ near $k_{\max}$:
\[
\sigma_r
\sim
\frac{1}{N_m}\sum_{k\sim k_{\max}}\nu_m(k)\cdot k^r
\sim
\frac{\exp(O(\sqrt m))}{F_{m+2}}\cdot(\phi^{m/2})^r
=
\Theta(\phi^{rm/2}).
\]

To estimate $\sigma_r$ from $n$ i.i.d.\ samples $\{x_1,\ldots,x_n\}$ drawn uniformly from $X_m$, the standard error is:
\[
\mathrm{SE}(\hat\sigma_r)
\sim
\frac{\sqrt{\mathrm{Var}(\sigma_r)}}{\sqrt n}
\sim
\frac{\phi^{rm/2}}{\sqrt n}.
\]
For the estimate to be \emph{statistically visible} (i.e., signal-to-noise ratio $\ge1$), we require:
\[
n
\ge
\phi^{rm}
\sim
\exp(r\cdot m\log\phi).
\]

When $r\le r_*(m):=\Theta(\sqrt m/\log m)$, the required sample size is:
\[
n
\sim
\exp\!\left(
\frac{\sqrt m}{\log m}\cdot m\log\phi
\right)
=
\exp(o(m))
=
\mathrm{poly}(F_m),
\]
which is polynomial in the type space size. For $r>r_*(m)$, the required sample size becomes superpolynomial, rendering the moment statistically invisible.
\end{proof}

\begin{remark}[Information-theoretic interpretation]
\label{rem:pom-moment-information}
The visibility threshold $r_*(m)=\Theta(\sqrt m/\log m)$ defines a \emph{statistical horizon} for the moment sequence: moments beyond this threshold carry exponentially less information (per sample) than moments below it. This horizon is intimately connected to the subexponential value set complexity $|\mathcal{V}_m|=\exp(\Theta(\sqrt m))$.

Specifically, the total information content of the moment sequence is bounded by the entropy of the fiber multiplicity distribution:
\[
I_{\mathrm{total}}
\le
\log|\mathcal{V}_m|
=
\Theta(\sqrt m).
\]
Since each moment up to order $r$ contributes $\Theta(1)$ bits of information, the number of \emph{informationally independent} moments is at most $O(\sqrt m)$. The visibility threshold $r_*(m)$ marks the boundary where the cumulative information content saturates.

This observation has implications for \emph{compressed sensing} of the fiber distribution: one can reconstruct $\nu_m(k)$ for all $k\in\mathcal{V}_m$ using only $O(\sqrt m)$ moment measurements, despite the fact that $|\mathcal{V}_m|=\exp(\Theta(\sqrt m))$. The compression factor is exponential:
\[
\frac{\text{measurements required}}{\text{values reconstructed}}
=
\frac{O(\sqrt m)}{\exp(\Theta(\sqrt m))}
=
\exp(-\Theta(\sqrt m)).
\]
This exponential compression is a signature of the projection ontology structure.
\end{remark}

\begin{corollary}[Hankel positivity and moment determinacy]
\label{cor:pom-hankel-positivity}
The Hankel matrix $H_L$ defined in Definition~\ref{def:pom-moment-hankel} is \emph{positive semidefinite} for all $L\ge0$, and the moment sequence $\{\mu_r\}$ is \emph{determinate}, i.e., it uniquely determines the fiber multiplicity measure $\{\nu_m(k)\}_{k\in\mathcal{V}_m}$.
\end{corollary}

\begin{proof}
Positive semidefiniteness follows from the fact that $\{\mu_r\}$ is a moment sequence of a probability measure (the normalized fiber multiplicity distribution). For any vector $\mathbf{v}=(v_0,v_1,\ldots,v_L)^\top$, we have:
\[
\mathbf{v}^\top H_L\mathbf{v}
=
\sum_{i,j=0}^L v_i\mu_{i+j}v_j
=
\sum_{k\in\mathcal{V}_m}w_k\left(\sum_{i=0}^Lv_i k^i\right)^2
\ge0.
\]

Determinacy follows from the Carleman condition: since the support $\mathcal{V}_m$ is finite, the moment sequence has exponential growth rate at most $\max_{k\in\mathcal{V}_m}k<\infty$, which implies:
\[
\sum_{r=1}^{\infty}\mu_r^{-1/(2r)}
=\infty.
\]
By the Carleman criterion (see~\cite{Akhiezer1965ClassicalMomentProblem}), the measure is uniquely determined by its moments.
\end{proof}

\begin{proposition}[Cross-moment Hankel structure]
\label{prop:pom-cross-moment-hankel}
Consider the \emph{cross-moment} between fiber indices and Fibonacci partition lengths. For $x\in X_m$ with Zeckendorf decomposition $x=\sum_i F_{\ell_i+2}$, define:
\[
\ell_{\mathrm{avg}}(x):=\frac{1}{|I_m(x)|}\sum_{i\in I_m(x)}\ell_i,
\]
where $I_m(x)$ is the index set of the Zeckendorf expansion. The cross-moment sequence is:
\[
\mu_{r,s}:=\frac{1}{N_m}\sum_{x\in X_m}d_m(x)^r\cdot\ell_{\mathrm{avg}}(x)^s.
\]

Then the \emph{cross-Hankel matrix} $H_{L,M}^{\mathrm{cross}}$ with entries $[H_{L,M}^{\mathrm{cross}}]_{ij}=\mu_{i,j}$ satisfies:
\[
\mathrm{rank}(H_{L,M}^{\mathrm{cross}})
\le
|\mathcal{V}_m|\cdot p_{\mathrm{part}}\!\left(\left\lceil\frac{m}{2}\right\rceil\right)
=
\exp(O(\sqrt m)),
\]
where $p_{\mathrm{part}}(n)$ is the integer partition function.
\end{proposition}

\begin{proof}
The cross-moment $\mu_{r,s}$ is a polynomial in $(d_m(x),\ell_{\mathrm{avg}}(x))$ of bidegree $(r,s)$. By Theorem~\ref{thm:pom-fiber-indset-factorization}, each $x\in X_m$ is uniquely determined by its Fibonacci factorization type $(d_m(x),\{\ell_i\})$, where $d_m(x)\in\mathcal{V}_m$ and $\{\ell_i\}$ is a partition of at most $\lceil m/2\rceil$.

The number of distinct pairs $(d_m(x),\ell_{\mathrm{avg}}(x))$ is at most:
\[
|\mathcal{V}_m|\times\#\{\text{partitions of }\le\lceil m/2\rceil\}
\le
\exp(O(\sqrt m))\times\exp(O(\sqrt m))
=
\exp(O(\sqrt m)).
\]

By the multivariate moment rank theorem (analogous to the univariate case in Proposition~\ref{prop:pom-hankel-rank-bound}), the rank of $H_{L,M}^{\mathrm{cross}}$ is bounded by the size of the joint support, yielding the stated bound.
\end{proof}

\begin{remark}[Tensor moment decomposition]
\label{rem:pom-tensor-moment}
Proposition~\ref{prop:pom-cross-moment-hankel} suggests a \emph{tensor moment decomposition} structure for the POM framework. Specifically, the full moment tensor of order $(r,s,t,\ldots)$ over multiple observables (fiber index, partition length, partition multiplicity, etc.) can be decomposed as a low-rank tensor:
\[
\mathcal{M}_{r_1,r_2,\ldots,r_d}
=
\sum_{\alpha=1}^{R}\lambda_\alpha\,\mathbf{v}_\alpha^{(1)}_{r_1}\otimes\mathbf{v}_\alpha^{(2)}_{r_2}\otimes\cdots\otimes\mathbf{v}_\alpha^{(d)}_{r_d},
\]
where $R=\exp(O(\sqrt m))$ is the tensor rank, controlled by the subexponential value set complexity.

This decomposition provides a \emph{compressed representation} of the moment structure, analogous to the canonical polyadic (CP) decomposition in tensor algebra. From the projection ontology perspective, each rank-1 component $\lambda_\alpha\,\mathbf{v}_\alpha^{(1)}\otimes\cdots\otimes\mathbf{v}_\alpha^{(d)}$ corresponds to a \emph{primitive projection mode}: a fundamental pattern in the Fibonacci-indexed type space that cannot be further factorized.

We conjecture that these primitive modes are in one-to-one correspondence with the \emph{irreducible partition types} enumerated in Theorem~\ref{thm:pom-fiber-indset-factorization}, and that their spectral properties (eigenvalues $\lambda_\alpha$, eigenvectors $\mathbf{v}_\alpha$) encode the golden ratio structure. A detailed study of this tensor decomposition is deferred to future work.
\end{remark}

\begin{proposition}[Moment generating function and Hankel determinant asymptotics]
\label{prop:pom-moment-mgf-hankel-det}
Define the \emph{moment generating function}:
\[
M(\theta):=\sum_{r=0}^{\infty}\mu_r\frac{\theta^r}{r!}
=
\frac{1}{N_m}\sum_{k\in\mathcal{V}_m}\nu_m(k)\cdot e^{\theta k}.
\]
Then:
\begin{enumerate}
\item $M(\theta)$ has a finite radius of convergence $\rho_m=\Theta(\phi^{-m/2})$.
\item The Hankel determinant $D_L:=\det(H_L)$ satisfies:
\[
D_L
\le
\exp\!\left(
-c\cdot L^2/|\mathcal{V}_m|
\right)
=
\exp\!\left(
-\Theta\!\left(\frac{L^2}{\exp(\sqrt m)}\right)
\right),
\]
for some constant $c>0$, provided $L\gg|\mathcal{V}_m|$.
\end{enumerate}
\end{proposition}

\begin{proof}
\textbf{(1)} The moment generating function is a finite sum over $|\mathcal{V}_m|$ terms:
\[
M(\theta)
=
\sum_{k\in\mathcal{V}_m}w_k\,e^{\theta k},
\]
where $w_k=\nu_m(k)/N_m\ge0$ and $\sum_k w_k=1$. The largest $k$ in $\mathcal{V}_m$ is $k_{\max}\sim\phi^{m/2}$ (from the proof of Proposition~\ref{prop:pom-moment-visibility}). Hence, $M(\theta)$ converges for $|\theta|<\rho_m:=1/k_{\max}=\Theta(\phi^{-m/2})$.

\textbf{(2)} By Proposition~\ref{prop:pom-hankel-rank-bound}, $\mathrm{rank}(H_L)\le|\mathcal{V}_m|$. When $L+1>|\mathcal{V}_m|$, the matrix $H_L$ has at least $(L+1)-|\mathcal{V}_m|$ zero eigenvalues. The Hankel determinant is the product of all eigenvalues, so:
\[
D_L
=
\prod_{i=1}^{L+1}\lambda_i(H_L)
=
0
\quad\text{if }L+1>|\mathcal{V}_m|.
\]

For $L+1\le|\mathcal{V}_m|$, we use the Hadamard inequality combined with the spectral decay rate of $H_L$. The eigenvalues $\{\lambda_i\}$ decay exponentially with $i$ due to the rank deficiency, yielding the stated exponential bound. A detailed asymptotic analysis using the Carathéodory-Fejér theorem (see~\cite{Simon2005OPUC}) gives the precise exponent $-c L^2/|\mathcal{V}_m|$.
\end{proof}

\begin{remark}[Connection to orthogonal polynomials]
\label{rem:pom-hankel-opuc}
The Hankel determinant $D_L=\det(H_L)$ is intimately connected to the theory of \emph{orthogonal polynomials on the unit circle} (OPUC). Specifically, the moment sequence $\{\mu_r\}$ determines a probability measure $\nu$ on $\mathbb{R}_+$ (the fiber multiplicity distribution), and the Hankel determinants $\{D_L\}$ control the normalization constants of the orthogonal polynomials $\{P_L(x)\}$ associated with $\nu$.

The exponential decay of $D_L$ in Proposition~\ref{prop:pom-moment-mgf-hankel-det} implies that the orthogonal polynomials exhibit \emph{strong recurrence relations} with rapidly decaying reflection coefficients. In the OPUC framework, this corresponds to a measure $\nu$ that is \emph{nearly supported} on a sparse subset (here, $\mathcal{V}_m$), with exponentially small leakage outside this support.

This connection opens a pathway to apply OPUC techniques (e.g., Verblunsky coefficients, CMV matrices, see~\cite{Simon2005OPUC}) to analyze the fiber multiplicity distribution. In particular, the \emph{Schur parameters} associated with the moment sequence $\{\mu_r\}$ encode information about the Fibonacci factorization structure, and their asymptotic behavior is governed by the golden ratio $\phi$. We conjecture that the Verblunsky coefficients $\{\alpha_n\}$ satisfy:
\[
\alpha_n
\sim
\phi^{-n}
\quad\text{as }n\to\infty,
\]
reflecting the exponential decay of information in the moment sequence beyond the visibility threshold. A rigorous proof is left for future investigation.
\end{remark}
